\documentclass[11pt]{letter}

\usepackage[utf8]{inputenc}
\usepackage[T1]{fontenc}
\usepackage[margin=1in]{geometry}
\usepackage{hyperref}
\usepackage{siunitx}

\signature{Jonathan Mace\\for the authors}
\address{Jonathan Mace\\Microsoft Research\\Redmond, WA, USA\\jonathanmace@microsoft.com}
\date{\today}

\begin{document}

\begin{letter}{%
  The Editor\\
  \textit{Inverse Problems in Science and Engineering}\\
  Taylor \& Francis}

\opening{Dear Editor,}

We respectfully submit the enclosed manuscript, ``Can You Hear the Shape of an
Organ? Practical Identifiability of Elastic Shell Parameters from Resonant
Frequencies,'' by Jonathan Mace and Brian R.\ Mace, for consideration in
\textit{Inverse Problems in Science and Engineering}.

The paper studies a practically important inverse question: what shell
parameters can be recovered from measured resonant frequencies, and when does
geometry make that recovery possible?  Using a fluid-filled oblate-shell Ritz
model together with its equivalent-sphere reduction as a degenerate comparator,
we show that the equivalent-sphere formulation is fundamentally
non-identifiable for spectral inversion, with
$\kappa \approx 1.37 \times 10^{10}$, whereas oblate eccentricity rescues
identifiability through mode-dependent curvature sampling, reducing the
condition number to $\kappa = 69.4$ at the canonical operating point.

We believe the manuscript fits \textit{Inverse Problems in Science and
Engineering} particularly well because its main contribution lies at the
interface of inverse-problems theory and engineering mechanics.  The paper is
not only about predicting shell resonances; it is about diagnosing when a
forward model that appears perfectly adequate for resonance prediction becomes
catastrophically ill-conditioned when used in reverse.  This
\emph{forward--inverse adequacy gap} is, in our view, the central result:
forward adequacy does not guarantee inverse adequacy.  That distinction should
be relevant to readers working on model updating, structural acoustics,
parameter estimation, and practical identifiability across applied inverse
problems.

The manuscript makes four linked contributions.  First, it establishes that the
equivalent-sphere reduction is not merely poorly conditioned but algebraically
rank-deficient for the inverse task.  Second, it shows that oblate geometry
lifts that degeneracy by diversifying modal curvature sampling across latitude.
Third, it organises the numerics into a three-class taxonomy of identifiability
regimes: singular, regular bounded, and non-lifting.  Fourth, it demonstrates a
practical conditioning floor, $\kappa_{\mathrm{floor}} \approx 269$, which sets
a noise-limited lower bound on achievable inversion quality within the present
Ritz framework.

More broadly, the manuscript connects a classical spectral question in inverse
problems to a concrete engineering setting: fluid-filled shell structures whose
frequencies may be easy to predict but difficult to invert reliably.  We hope
this combination of conceptual clarity and application relevance will make the
paper of interest to the journal's readership.

The manuscript is original, has not been published previously, is not under
consideration elsewhere, and has been approved by both authors.  The authors
declare no competing financial interests or personal relationships that could
have influenced the work.  All analysis code and manuscript sources are
available in the accompanying repository for reproducibility.

Thank you for your consideration.

\closing{Yours sincerely,}

\end{letter}
\end{document}
